% !TEX root = ../main.tex

\chapter{Introduction}

\section{Présentation personnelle}

Actuellement étudiant en [intitulé de la formation], j'ai effectué un
stage de [durée] au sein de l'entreprise ISI.nc, du [date de début] au
[date de fin]. Mon parcours m'a progressivement orienté vers le
développement web, et plus particulièrement vers le développement full
stack, à la croisée du frontend et du backend. Ce stage représentait pour
moi l'opportunité de mettre en pratique mes compétences sur un projet réel,
de taille significative, au contact d'une équipe expérimentée.

J'ai choisi ce stage car [motivation personnelle : intérêt pour la
mission proposée, réputation de l'entreprise, technologies utilisées,
etc. — à personnaliser].

\section{Présentation de l'entreprise d'accueil}

ISI.nc est une entreprise spécialisée dans le développement de solutions
logicielles sur mesure. Elle intervient auprès de divers clients pour
concevoir, développer et maintenir des applications métier adaptées à
leurs besoins spécifiques. L'entreprise s'appuie sur des méthodologies
agiles et une expertise technique couvrant l'ensemble de la chaîne de
développement, du frontend au backend, en passant par l'infrastructure.

\section{Présentation du client FPG et du contexte métier}

Le stage s'est déroulé dans le cadre d'une mission réalisée pour le
compte du FPG (Fonds Paritaire de Gestion, ou équivalent selon le
contexte réel), un organisme chargé de la gestion des demandes de
financement de formation professionnelle déposées par des structures
adhérentes (entreprises).

L'application développée accompagne l'ensemble du cycle de vie d'une
demande de financement :

\begin{itemize}
  \item le \textbf{dépôt} de la demande par la structure adhérente ;
  \item l'\textbf{instruction} du dossier par les équipes du FPG ;
  \item l'\textbf{engagement} financier associé à la demande ;
  \item le \textbf{règlement} des sommes dues.
\end{itemize}

Ce cycle mobilise plusieurs types d'entités métier interconnectées :
stagiaires, montages financiers, tiers bénéficiaires, règlements,
documents justificatifs, structures adhérentes et comptes collaborateurs.
La complexité et la sensibilité de ces données — notamment financières —
rendent la traçabilité des modifications apportées à ces dossiers
particulièrement critique, ce qui constitue précisément le cœur de la
mission qui m'a été confiée.

\section{Annonce du plan}

Ce rapport présente dans un premier temps le contexte de travail et la
méthodologie adoptée, ainsi que la stack technique du projet. Il détaille
ensuite en profondeur la mission principale du stage : la conception et
le développement d'un système d'historisation des modifications. Les
compétences acquises et les difficultés rencontrées sont ensuite
discutées, avant de conclure sur le bilan global de cette expérience.
